
% =========================================================
% Driver Text: Gallian, Contemporary Abstract Algebra (7th ed.)
% File: abstract-algebra/abstract-algebra.tex
% =========================================================

\section{Abstract Algebra}

% =========================================================
% =========================================================

% =========================================================
% I. Integers and Equivalence
% =========================================================

\begin{enumerate}

\item \textbf{Foundations: Integers and Equivalence Relations}

\begin{enumerate}
    \item Modular arithmetic
    \item Equivalence relations and partitions
    \item Functions and mappings
\end{enumerate}

\vspace{0.5em}
\textbf{Structural Theme:}
Congruence and equivalence encode algebraic structure via partitions.
This stage formalizes symmetry at the level of arithmetic.

% =========================================================
% II. Group Theory
% =========================================================

\item \textbf{Groups}

\begin{enumerate}
    \item Definition and basic properties
    \item Subgroups
    \item Cyclic groups
    \item Permutation groups
    \item Isomorphisms
    \item Cosets and Lagrange's Theorem
    \item Normal subgroups and factor groups
    \item Homomorphisms
    \item Fundamental Theorem of Finite Abelian Groups
\end{enumerate}

\vspace{0.5em}
\textbf{Structural Theme:}
Groups formalize symmetry.
Cosets measure deviation from subgroup structure.
Normality permits quotient construction.
Homomorphisms reveal structural preservation.
Classification emerges in the finite abelian case.

% =========================================================
% III. Ring Theory
% =========================================================

\item \textbf{Rings}

\begin{enumerate}
    \item Definition and examples
    \item Integral domains
    \item Ideals and factor rings
    \item Ring homomorphisms
    \item Polynomial rings
    \item Unique factorization domains
    \item Euclidean domains
\end{enumerate}

\vspace{0.5em}
\textbf{Structural Theme:}
Rings generalize arithmetic.
Ideals control structure via quotients.
Factorization encodes algebraic rigidity.
Euclidean domains permit algorithmic structure.

% =========================================================
% IV. Field Theory
% =========================================================

\item \textbf{Fields and Extensions}

\begin{enumerate}
    \item Vector spaces
    \item Field extensions
    \item Algebraic extensions
    \item Finite fields
    \item Geometric constructions
\end{enumerate}

\vspace{0.5em}
\textbf{Structural Theme:}
Fields unify arithmetic and linear algebra.
Extensions enlarge solvability.
Finite fields exhibit deep classification.
Constructibility links algebra and geometry.

% =========================================================
% V. Advanced Topics
% =========================================================

\item \textbf{Special Topics}

\begin{enumerate}
    \item Sylow Theorems
    \item Finite simple groups
    \item Generators and relations
    \item Symmetry groups
    \item Group actions and counting (Burnside)
    \item Coding theory
    \item Galois theory
    \item Cyclotomic extensions
\end{enumerate}

\vspace{0.5em}
\textbf{Structural Theme:}
Group actions connect algebra to combinatorics.
Sylow theory controls finite structure.
Galois theory connects symmetry and solvability.
Cyclotomic extensions bridge algebra and number theory.

\end{enumerate}

% =========================================================
% Structural Remarks
% =========================================================

\vspace{1em}

\begin{remark}
Abstract algebra studies algebraic structures defined by operations
subject to axioms. Each structure (group, ring, field) is a
package consisting of:
\[
\text{Set} + \text{Operations} + \text{Axioms}.
\]
\end{remark}

\begin{remark}
The central construction principle is:
\[
\textbf{Substructure} \longrightarrow \textbf{Quotient Structure}.
\]
Normal subgroups and ideals make quotient constructions possible.
\end{remark}

\begin{remark}[Structural Position]
Abstract algebra in this project builds on:

\begin{itemize}
    \item Set theory and logic (Phase 0)
    \item Linear algebra (vector space structure)
    \item Real analysis (structural proof discipline)
\end{itemize}

Group theory forms the conceptual backbone.
Ring theory extends arithmetic.
Field theory culminates in symmetry of roots (Galois theory).
\end{remark}

% =========================================================
% Input Files (To Be Developed)
% =========================================================

% -----------------------
% Foundations
% -----------------------
\section{Integers}
% =========================================================
% NOTES TEMPLATE (Definitions $\to$ Theorems/Proofs $\to$ Consequences)
% File: notes-integers.tex
% =========================================================

\subsection{Integers}
\label{sec:integers}  


% ---------------------------------------------------------
\subsection*{Basic Definitions and Theorems}


% =========================================================
% NOTES TEMPLATE (Definitions $\to$ Theorems/Proofs $\to$ Consequences)
% File: notes-modular-arithmetic.tex
% =========================================================

\section{Modular Arithmetic}
\label{sec:modular-arithmetic}

% ---------------------------------------------------------
\subsection*{Basic Definitions and Theorems}


% =========================================================
% Induction on the integers
% =========================================================

\subsection{Induction}

% ---------------------------------------------------------
\subsection*{Definitions and Theorems}
% ---------------------------------------------------------





















% -----------------------
% Groups
% -----------------------
% \input{volume-vi/book-algebra/abstract-algebra/notes/notes-groups}
% \input{volume-vi/book-algebra/abstract-algebra/notes/notes-subgroups}
% \input{volume-vi/book-algebra/abstract-algebra/notes/notes-cyclic-groups}
% \input{volume-vi/book-algebra/abstract-algebra/notes/notes-permutation-groups}
% \input{volume-vi/book-algebra/abstract-algebra/notes/notes-isomorphisms}
% \input{volume-vi/book-algebra/abstract-algebra/notes/notes-cosets}
% \input{volume-vi/book-algebra/abstract-algebra/notes/notes-normal-subgroups}
% \input{volume-vi/book-algebra/abstract-algebra/notes/notes-homomorphisms}
% \input{volume-vi/book-algebra/abstract-algebra/notes/notes-finite-abelian-groups}

% -----------------------
% Rings
% -----------------------
% \input{volume-vi/book-algebra/abstract-algebra/notes/notes-rings}
% \input{volume-vi/book-algebra/abstract-algebra/notes/notes-ideals}
% % =========================================================
% Polynomial Rings
% =========================================================

% ---------------------------------------------------------
\section{Preliminary Definitions}
% ---------------------------------------------------------

\begin{definition}[Ring]
\label{def:algebraic-geometry-ring}
A set $R$ with two binary operations $+ : R \times R \to R$
and $\cdot : R \times R \to R$ is a \emph{ring} if:

\begin{enumerate}
\item \textbf{Additive structure:}
$(R, +)$ is an abelian group.

\item \textbf{Associativity of multiplication:}
\[
(a \cdot b) \cdot c = a \cdot (b \cdot c)
\quad \text{for all } a,b,c \in R.
\]

\item \textbf{Multiplicative identity:}
There exists $1 \in R$ such that
\[
1 \cdot a = a \cdot 1 = a
\quad \text{for all } a \in R.
\]

\item \textbf{Distributivity:}
\[
a \cdot (b + c) = a \cdot b + a \cdot c
\quad \text{and} \quad
(a + b) \cdot c = a \cdot c + b \cdot c
\quad \text{for all } a,b,c \in R.
\]
\end{enumerate}
\end{definition}

\begin{remark*}[Standard quantified statement]
A ring is a set equipped with addition and multiplication satisfying the
additive-group, multiplicative-associativity, multiplicative-identity, and
distributive laws.
\end{remark*}

\begin{remark*}[Interpretation]
A ring generalizes a field by dropping the requirement that
nonzero elements have multiplicative inverses,
and by not requiring multiplication to be commutative.
When multiplication is commutative, we call $R$ a
\emph{commutative ring}.
\end{remark*}

\begin{definition}[Commutative Ring]
\label{def:algebraic-geometry-commutative-ring}
A ring $R$ is \emph{commutative} if
\[
a \cdot b = b \cdot a
\quad \text{for all } a, b \in R.
\]
\end{definition}

\begin{remark*}[Standard quantified statement]
A ring $R$ is commutative when multiplication in $R$ is commutative.
\end{remark*}

\begin{remark*}[Structural Summary]
The hierarchy of algebraic structures encountered so far is:

\begin{center}
\begin{tabular}{l|c|c}
Structure & Additive Group & Multiplicative Structure \\
\hline
Ring            & Abelian & Associative, with identity \\
Commutative Ring & Abelian & Associative, commutative, with identity \\
Field           & Abelian & Abelian on nonzero elements
\end{tabular}
\end{center}

Every field is a commutative ring, but not every commutative ring is a field.
\end{remark*}

% ---------------------------------------------------------
\subsection*{Polynomial Rings}
% ---------------------------------------------------------

\begin{definition}[Polynomial]
\label{def:polynomial-over-ring}
Let $R$ be a commutative ring.
A \emph{polynomial} in one variable over $R$ is a formal expression
\[
f = a_n x^n + a_{n-1} x^{n-1} + \cdots + a_1 x + a_0,
\]
where $n \in \mathbb{N}$, the \emph{coefficients} $a_0, a_1, \dots, a_n \in R$,
and $x$ is a formal symbol called an \emph{indeterminate}.
\end{definition}

\begin{remark*}[Standard quantified statement]
A polynomial over $R$ is a finite formal sum $\sum_i a_i x^i$ with
coefficients $a_i \in R$.
\end{remark*}

\begin{remark*}[Interpretation]
The indeterminate $x$ is not a variable ranging over values in $R$.
It is a formal placeholder that encodes the coefficient sequence.
Two polynomials are equal if and only if all their coefficients are equal.
\end{remark*}

\begin{definition}[Degree]
\label{def:degree-of-polynomial}
Let $f = a_n x^n + \cdots + a_0$ be a polynomial over $R$.
If $a_n \neq 0$, then the \emph{degree} of $f$ is
\[
\deg(f) := n.
\]
The coefficient $a_n$ is called the \emph{leading coefficient} of $f$.
A polynomial with leading coefficient $1$ is called \emph{monic}.
\end{definition}

\begin{remark*}[Standard quantified statement]
For a nonzero polynomial $f=a_nx^n+\cdots+a_0$ with leading coefficient
$a_n\neq 0$, the degree of $f$ is $n$.
\end{remark*}

\begin{remark*}[Interpretation]
The zero polynomial $f = 0$ has no leading coefficient.
Its degree is left undefined, or assigned $-\infty$ by convention
to preserve the identity $\deg(fg) = \deg(f) + \deg(g)$.
\end{remark*}

\begin{definition}[Polynomial Ring]
\label{def:polynomial-ring}
Let $R$ be a commutative ring.
The \emph{polynomial ring} over $R$ in one indeterminate $x$,
denoted $R[x]$, is the set of all polynomials in $x$ with
coefficients in $R$, equipped with addition and multiplication
defined by:
\[
\left(\sum_{i} a_i x^i\right) + \left(\sum_{i} b_i x^i\right)
:= \sum_{i} (a_i + b_i) x^i,
\]
\[
\left(\sum_{i} a_i x^i\right) \cdot \left(\sum_{j} b_j x^j\right)
:= \sum_{k} \left(\sum_{i+j=k} a_i b_j\right) x^k.
\]
\end{definition}

\begin{remark*}[Standard quantified statement]
The set $R[x]$ is the collection of finite formal sums in $x$ with
coefficients in $R$, equipped with coefficientwise addition and Cauchy-product
multiplication.
\end{remark*}

\begin{example*}[Polynomial arithmetic over the integers]
Let $R = \mathbb{Z}$ and consider
\[
f = 2x^2 + 3x + 1, \qquad g = x + 4 \quad \in \mathbb{Z}[x].
\]
Then:
\[
f + g = 2x^2 + 4x + 5,
\]
\[
f \cdot g = 2x^3 + 11x^2 + 13x + 4.
\]
\end{example*}

\begin{remark*}[Structural Summary]
With these operations, $R[x]$ is itself a commutative ring.
The original ring $R$ embeds into $R[x]$ as the constant polynomials.

\begin{center}
\begin{tabular}{l|c}
Property & Holds in $R[x]$? \\
\hline
Commutative ring & Always \\
Field & Only in degenerate cases \\
$R$ embeds in $R[x]$ & Always
\end{tabular}
\end{center}
\end{remark*}

% ---------------------------------------------------------
\subsection*{Polynomial Rings in Several Variables}
% ---------------------------------------------------------

\begin{definition}[Polynomial Ring in $n$ Variables]
\label{def:polynomial-ring-in-n-variables}
Let $R$ be a commutative ring and let $n \in \mathbb{N}$.
The \emph{polynomial ring in $n$ variables} over $R$, denoted
\[
R[x_1, x_2, \dots, x_n],
\]
is defined inductively by
\[
R[x_1, \dots, x_n] := R[x_1, \dots, x_{n-1}][x_n].
\]
Its elements are finite sums of the form
\[
f = \sum_{\alpha} a_\alpha \, x^\alpha,
\]
where the sum ranges over multi-indices
$\alpha = (\alpha_1, \dots, \alpha_n) \in \mathbb{N}^n$,
the coefficients $a_\alpha \in R$, and
$x^\alpha := x_1^{\alpha_1} \cdots x_n^{\alpha_n}$.
\end{definition}

\begin{remark*}[Standard quantified statement]
The ring $R[x_1,\dots,x_n]$ consists of finite formal sums
$\sum_\alpha a_\alpha x^\alpha$ indexed by multi-indices.
\end{remark*}

\begin{definition}[Monomial]
\label{def:monomial}
A \emph{monomial} in $R[x_1, \dots, x_n]$ is a polynomial
of the form
\[
x^\alpha = x_1^{\alpha_1} \cdots x_n^{\alpha_n}
\]
for some multi-index $\alpha \in \mathbb{N}^n$.
\end{definition}

\begin{remark*}[Standard quantified statement]
A monomial in $R[x_1,\dots,x_n]$ is a single multi-index power $x^\alpha$.
\end{remark*}

\begin{definition}[Total Degree]
\label{def:total-degree}
The \emph{total degree} of the monomial $x^\alpha$ is
\[
|\alpha| := \alpha_1 + \alpha_2 + \cdots + \alpha_n.
\]
The degree of a polynomial $f \in R[x_1,\dots,x_n]$ is
the maximum total degree among all monomials with nonzero coefficient.
\end{definition}

\begin{remark*}[Standard quantified statement]
The total degree of $x^\alpha$ is $|\alpha|=\alpha_1+\cdots+\alpha_n$.
\end{remark*}

\begin{example*}[Example]
In $\mathbb{R}[x,y,z]$, the polynomial
\[
f = 3x^2 y + xy^2 z - 5z^3
\]
has three terms with total degrees $3$, $4$, and $3$ respectively.
Hence $\deg(f) = 4$.
\end{example*}

\begin{remark*}[Relevance to Algebraic Geometry]
Polynomial rings in several variables are the foundational
algebraic object of algebraic geometry.
The geometric objects studied in subsequent sections ---
affine varieties, ideals, coordinate rings ---
are all defined in terms of $k[x_1, \dots, x_n]$
for a field $k$.
The interplay between the algebra of $k[x_1,\dots,x_n]$
and the geometry of its zero sets is the central theme
of Clader--Ross.
\end{remark*}

% ---------------------------------------------------------
\subsection*{Ideals}
% ---------------------------------------------------------

\begin{definition}[Ideal]
\label{def:ideal}
Let $R$ be a commutative ring.
A subset $I \subseteq R$ is an \emph{ideal} of $R$ if:
\begin{enumerate}
\item $0 \in I$,
\item $a + b \in I$ for all $a, b \in I$,
\item $r \cdot a \in I$ for all $r \in R$ and $a \in I$.
\end{enumerate}
\end{definition}

\begin{remark*}[Standard quantified statement]
An ideal of $R$ is an additive subgroup of $R$ closed under multiplication by
arbitrary elements of $R$.
\end{remark*}

\begin{definition}[Generated Ideal]
\label{def:generated-ideal}
Let $R$ be a commutative ring and let $f_1, \dots, f_m \in R$.
The \emph{ideal generated by} $f_1, \dots, f_m$ is
\[
\langle f_1, \dots, f_m \rangle
:= \left\{ \sum_{i=1}^m r_i f_i : r_i \in R \right\}.
\]
An ideal of this form is called \emph{finitely generated}.
\end{definition}

\begin{remark*}[Standard quantified statement]
The ideal generated by $f_1,\dots,f_m$ is the set of all finite
$R$-linear combinations of those generators.
\end{remark*}

\begin{definition}[Finitely Generated Ideal]
\label{def:finitely-generated-ideal}
An ideal $I \subseteq R$ is \emph{finitely generated} if there exist
$f_1, \dots, f_m \in R$ such that $I = \langle f_1, \dots, f_m \rangle$.
\end{definition}

\begin{remark*}[Standard quantified statement]
An ideal $I$ is finitely generated when $I=\langle f_1,\dots,f_m\rangle$
for some finite list $f_1,\dots,f_m$ of elements of $R$.
\end{remark*}

\begin{definition}[Noetherian Ring]
\label{def:noetherian-ring}
A commutative ring $R$ is \emph{Noetherian} if every ideal of $R$
is finitely generated.
\end{definition}

\begin{remark*}[Standard quantified statement]
A commutative ring $R$ is Noetherian when every ideal $I\subseteq R$ is
finitely generated.
\end{remark*}

\begin{theorem}[Hilbert Basis Theorem]
\label{thm:hilbert-basis-theorem}
If $R$ is a Noetherian ring, then the polynomial ring $R[x]$
is also Noetherian.

In particular, $k[x_1, \dots, x_n]$ is Noetherian for any field $k$.
\hyperref[prf:hilbert-basis-theorem]{Go to proof.}
\end{theorem}

\begin{remark*}[Standard quantified statement]
If $R$ is Noetherian, then $R[x]$ is Noetherian.
\end{remark*}

\begin{remark*}[Interpretation]
The Hilbert Basis Theorem guarantees that every ideal in
$k[x_1, \dots, x_n]$ is finitely generated.
This is a foundational finiteness result:
it means every algebraic variety can be cut out by
finitely many polynomial equations.
\end{remark*}

% ---------------------------------------------------------
\subsection*{Quotient Rings}
% ---------------------------------------------------------

\begin{definition}[Quotient Ring]
\label{def:quotient-ring}
Let $R$ be a commutative ring and let $I \subseteq R$ be an ideal.
The \emph{quotient ring} $R/I$ is the set of cosets
\[
R/I := \{ a + I : a \in R \},
\]
equipped with addition and multiplication defined by
\[
(a + I) + (b + I) := (a + b) + I,
\]
\[
(a + I) \cdot (b + I) := (a \cdot b) + I.
\]
\end{definition}

\begin{remark*}[Standard quantified statement]
For an ideal $I$ of $R$, the quotient $R/I$ is the set of additive cosets of
$I$ with operations induced from $R$.
\end{remark*}

\begin{remark*}[Interpretation]
The quotient ring $R/I$ is a commutative ring.
Elements of $R/I$ are equivalence classes under the relation
$a \sim b \iff a - b \in I$.
\end{remark*}

\begin{example*}[A complex quotient]
In $\mathbb{R}[x]$, consider the ideal $I = \langle x^2 + 1 \rangle$.
The quotient ring
\[
\mathbb{R}[x] / \langle x^2 + 1 \rangle \;\cong\; \mathbb{C},
\]
since in the quotient, $x$ satisfies $x^2 = -1$,
which is precisely the defining relation of $i \in \mathbb{C}$.
\end{example*}

\begin{remark*}[Structural Position]
Quotient rings of polynomial rings are the coordinate rings
of affine varieties, which are studied in depth in the next chapter.
The passage
\[
k[x_1,\dots,x_n] \;\longrightarrow\; k[x_1,\dots,x_n]/I
\]
is the algebraic encoding of restricting from all of $\mathbb{A}^n$
to the variety $V(I)$.
\end{remark*}

% \input{volume-vi/book-algebra/abstract-algebra/notes/notes-ufd-euclidean}

% -----------------------
% Fields
% -----------------------
% \input{volume-vi/book-algebra/abstract-algebra/notes/notes-field-extensions}
% \input{volume-vi/book-algebra/abstract-algebra/notes/notes-finite-fields}
% \input{volume-vi/book-algebra/abstract-algebra/notes/notes-galois-theory}

% -----------------------
% Special Topics
% -----------------------
% \input{volume-vi/book-algebra/abstract-algebra/notes/notes-sylow}
% \input{volume-vi/book-algebra/abstract-algebra/notes/notes-group-actions}
% \input{volume-vi/book-algebra/abstract-algebra/notes/notes-coding-theory}
